\documentclass[11pt]{article}

\usepackage[margin=1in]{geometry}
\usepackage{booktabs}
\usepackage{array}
\usepackage{hyperref}
\usepackage{natbib}
\usepackage{xcolor}
\usepackage{xspace}

\title{Verifier-Calibrated Search and Repair for Faithful Text-to-PDDL Generation}
\author{Pritam Mondal}
\date{}

\newcommand{\vcsr}{\textsc{VCSR}\xspace}
\newcommand{\pddl}{\textsc{PDDL}\xspace}

\begin{document}
\maketitle

\begin{abstract}
Natural-language interfaces to automated planning promise to let users specify
planning problems without manually writing Planning Domain Definition Language
(\pddl) artifacts. However, generated \pddl can be syntactically valid, and even
solvable, while still formalizing the wrong task. This is one instance of a
broader problem: natural-language systems can produce structured outputs that
are valid in form but unfaithful in meaning. We study this valid-but-wrong
failure mode through semantic equivalence on Planetarium. We introduce \vcsr,
a verifier-calibrated search-and-repair framework for faithful text-to-\pddl
generation. \vcsr generates a small candidate pool, ranks
parseable candidates with a learned semantic verifier, and applies one
domain-aware repair call to the verifier-selected candidate. On an in-domain
Planetarium evaluation over \texttt{blocksworld} and \texttt{gripper}, using
untouched final seeds, full \vcsr improves mean semantic equivalence at $K=8$
from 0.368 for prompt-only generation to 0.772, a gain of 40.4 percentage
points. Verifier-only search improves to 0.420, making it a useful scaffold but
not the final solution. The original $K=8$ candidate-pool oracle is only 0.464;
\vcsr exceeds it because repair creates a new candidate outside the original
pool. We also report negative evidence: planner/solvability search and further
verifier training did not reliably solve semantic selection. The final system is
a large net win, especially on \texttt{gripper}, but repair can hurt
already-correct \texttt{blocksworld} candidates.
\end{abstract}

\section{Introduction}

Automated planning systems rely on formal models of states, actions, and goals.
The Planning Domain Definition Language (\pddl) has long served as a standard
representation for such models \citep{ghallab1998pddl,fox2003pddl21}. If users
could describe a task in natural language and reliably obtain faithful \pddl,
classical planners would become substantially easier to deploy. Large language
models make this goal tempting because they can produce structured, code-like
artifacts from short descriptions. The danger is that a valid-looking artifact
can still be the wrong formalization.

Planetarium makes this gap concrete. It evaluates text-to-\pddl translation by
semantic equivalence rather than merely by parsing or plan existence
\citep{zuo2024planetarium}. This distinction matters: a generated problem can
parse and admit a plan while changing objects, initial facts, or goals. A
planner may then solve the formal problem perfectly while solving the wrong
user task.

\paragraph{A simple example.}
Imagine asking an assistant to turn ``deliver the package to Purbita after Utpal
signs for it'' into a checklist. A checklist that says ``Utpal signs after Purbita
receives the package'' may be perfectly well formatted, but it encodes the wrong
task. Text-to-\pddl has the same problem in a stricter symbolic setting: a file
can parse and even be solvable while reversing, dropping, or subtly changing
the user's intended constraints. \vcsr targets this broader problem of
faithful formalization, with Planetarium providing a rigorous testbed.

Prior work has explored LLM-driven planning formalization, including NL2Plan
\citep{gestrin2024nl2plan}, planning-domain generation
\citep{oswald2024planningdomains}, and studies of the limits of LLMs as
planning formalizers \citep{huang2024planningformalizers}. These systems show
that language models can be useful planning front ends, but they also reinforce
the need for mechanisms that target semantic faithfulness rather than only
syntactic validity or solvability.

We propose \vcsr: Verifier-Calibrated Search and Repair. \vcsr is built around
the view that faithful text-to-\pddl generation requires a system, not only a
larger prompt or a better verifier. The system should generate alternatives,
score semantic plausibility, and repair promising but wrong formalizations.
Figure~\ref{fig:pipeline} summarizes the pipeline.

\begin{figure}[t]
\centering
\fbox{
\begin{minipage}{0.93\linewidth}
\centering
Natural-language task
$\rightarrow$ generate $K$ \pddl candidates
$\rightarrow$ parse candidates
$\rightarrow$ score with frozen semantic verifier
$\rightarrow$ select verifier-ranked candidate
$\rightarrow$ one-step domain-aware repair
$\rightarrow$ evaluate semantic equivalence.
\end{minipage}}
\caption{\vcsr pipeline. The repair prompt receives the selected candidate and
non-oracle feedback, but never gold \pddl or an equivalence label.}
\label{fig:pipeline}
\end{figure}

This paper makes four contributions. First, we present a search-and-repair
framework for faithful text-to-\pddl generation. Second, we use fixed-pool
replay to separate candidate-pool randomness from verifier and policy changes.
Third, we show on untouched final seeds that repair-augmented \vcsr improves
semantic equivalence over prompt-only generation, random parseable best-of-$K$,
planner/solvability search, and verifier-only search. Fourth, we report
negative results that clarify why the final gain comes from repair: additional
verifier training and simple solvability-based search were not enough.

\section{Related Work}

\paragraph{Text-to-\pddl and planning formalization.}
Planetarium is the closest benchmark to our setting. It provides natural
language descriptions paired with \pddl problems and evaluates generated
problems using equivalence-based semantics \citep{zuo2024planetarium}. NL2Plan
studies robust LLM-driven planning from minimal text descriptions
\citep{gestrin2024nl2plan}. Other recent work investigates LLM-generated
planning domains \citep{oswald2024planningdomains}, generalized planning in
\pddl domains with pretrained language models
\citep{silver2024generalizedplanning}, and the degradation of LLM
formalization under more natural descriptions
\citep{huang2024planningformalizers}. Our work differs in emphasizing semantic
equivalence of generated problem files and in making repair a first-class
system stage.

\paragraph{Best-of-$K$ generation and test-time search.}
Sampling multiple outputs and selecting among them is a common way to improve
language-model reasoning and generation. Self-consistency demonstrated the
value of sampling diverse reasoning paths and aggregating outcomes
\citep{wang2023selfconsistency}. More directly related to symbolic artifacts,
recent work studies test-time scaling for generating symbolic world models
\citep{yu2025symbolicworldmodels}. \vcsr follows the same broad test-time
compute principle, but our final system does not only select from the original
candidate pool; it repairs the verifier-selected candidate.

\paragraph{Repair and refinement.}
LLM refinement methods such as Self-Refine use feedback to iteratively improve
model outputs \citep{madaan2023selfrefine}. \vcsr uses a narrower and more
controlled repair stage: a single repair call is applied after verifier-guided
selection, and the repair prompt receives no gold \pddl or equivalence label.
This keeps the repair stage close to a deployable setting where the system can
inspect its own candidate and lightweight diagnostics but does not know the
ground truth.

\paragraph{Planning tools and weak correctness signals.}
Classical planners and validators remain important planning tools. Fast
Downward is a widely used planning system \citep{helmert2006fastdownward}, and
VAL supports plan validation for \pddl \citep{howey2004val}. In this paper,
however, solvability and validation are diagnostic signals rather than final
correctness metrics. The primary metric is semantic equivalence to the intended
task.

\paragraph{Calibration and selective prediction.}
During development we studied verifier calibration and risk-coverage, motivated
by the miscalibration of neural confidence scores \citep{guo2017calibration}
and selective prediction work on risk-coverage tradeoffs
\citep{geifman2017selective}. The final paper-facing mechanism is not
abstention, however. We use the frozen verifier as a search scaffold and repair
as the decisive system-level intervention.

\section{Method}

\subsection{Task}

We study text-to-\pddl problem generation in fixed in-domain Planetarium tasks.
Each example consists of a natural-language description and a gold \pddl
problem. A generated candidate is correct if it is semantically equivalent to
the gold task under Planetarium's equivalence evaluator
\citep{zuo2024planetarium,planetarium_dataset,planetarium_code}. We evaluate on
the \texttt{test} split and restrict the final claim to \texttt{blocksworld} and
\texttt{gripper}. We make no final claim for \texttt{floor-tile}.

\subsection{Candidate Generation}

Given an input description $x$, the generator produces $K$ candidate \pddl
problems $\{y_1,\ldots,y_K\}$ using stochastic decoding. Our main experiments
use $K \in \{1,4,8\}$, with $K=8$ as the primary operating point. Candidate
outputs are parsed and normalized by the project evaluation stack. The $K=1$
condition is the prompt-only baseline: the model is asked once, and the first
candidate is evaluated directly.

\subsection{Semantic Verifier}

The verifier is a DeBERTaV3-base cross-encoder \citep{he2021debertav3}. It
scores pairs $(x,y)$, where $x$ is the natural-language task and $y$ is a
candidate \pddl problem. It was trained from automatically labeled examples
derived from Planetarium equivalence: parseable equivalent candidates are
positive examples, while parseable non-equivalent candidates and targeted
perturbations are negative examples. The final paper-facing verifier is the
round-4 checkpoint recorded in
\path{results/verifier/best_current/selection.yaml}. Later verifier training
experiments are treated as diagnostics and are not promoted.

\subsection{Selection Policies}

We evaluate four policies. \texttt{greedy\_first} selects the first generated
candidate. This is the prompt-only policy at $K=1$ and a simple baseline at
larger $K$. \texttt{random\_parseable} selects uniformly from parseable
candidates. \texttt{verifier\_ranked} selects the parseable candidate with the
highest frozen-verifier score. \texttt{verifier\_ranked\_repair} first selects
the verifier-ranked candidate at $K=8$ and then applies one repair call. None
of these policies uses the gold \pddl or equivalence label at selection time.

\subsection{One-Step Domain-Aware Repair}

The final \vcsr policy applies one repair call to the verifier-selected $K=8$
candidate. The repair prompt receives the natural-language task, the selected
candidate \pddl, and lightweight non-oracle feedback. It never receives the
gold \pddl, the equivalence label, or an oracle-best candidate. For
\texttt{blocksworld}, we use the generic repair prompt. For \texttt{gripper},
we use a Planetarium-specific repair prompt that encourages the object and
predicate conventions needed by the benchmark. This domain-aware repair policy
is the only promoted repair variant. A later score-guarded repair variant is
reported as a negative follow-up because it did not reduce blocksworld hurts.

\section{Experimental Setup}

\paragraph{Dataset and final gate.}
The final evaluation uses Planetarium \texttt{test} rows from
\texttt{blocksworld} and \texttt{gripper}. The primary final gate uses untouched
seeds 51--55, 50 rows per seed, for 250 total rows. These seeds were not used
for prompt tuning, verifier checkpoint selection, repair-prompt tuning, or
guard selection. We evaluate $K \in \{1,4,8\}$ and use the frozen round-4
verifier from \path{results/verifier/best_current/selection.yaml}.

\paragraph{Metrics.}
The primary metric is semantic equivalence rate. We also report parse rate,
equivalence given parse, domain and style slices, candidate-pool oracle
diagnostics, and repair helped/hurt/tied counts. ``Helped'' means the original
verifier-selected candidate was non-equivalent and the repaired candidate was
equivalent. ``Hurt'' means the original was equivalent and repair made it
non-equivalent. ``Tied'' includes both-success and both-fail cases.

\paragraph{Oracle best-of-$K$ diagnostic.}
Oracle best-of-$K$ is an offline candidate-pool diagnostic, not a deployable
policy. It uses ground-truth equivalence labels after generation to ask whether
at least one of the original $K$ candidates is semantically equivalent. It is
therefore an upper bound for policies that must choose among the original
candidates. It is not an upper bound for repair-augmented \vcsr, because repair
creates a new candidate outside the original pool.

\paragraph{Replay and reproducibility.}
During development, we used cached candidate pools to replay alternative
verifiers and policies on identical candidates. This avoided confounding policy
changes with generator randomness and was used to reject additional verifier
training rounds and planner/solvability-only search policies. The final
evidence is stored under
\path{results/vcsr/final_repair_gate_round4}; derived paper tables are stored
under \path{results/paper/final_vcsr}.

\section{Results}

\subsection{Main Final Gate}

Table~\ref{tab:main} shows the final aggregate result over untouched seeds
51--55. The primary operating point is $K=8$. Prompt-only generation reaches
0.368 semantic equivalence. Full repair-augmented \vcsr reaches 0.772, which
means 77.2\% of rows are semantically correct. The absolute improvement over
prompt-only is +0.404, or +40.4 percentage points. Verifier-only search reaches
0.420, confirming that the verifier has useful signal, but the large jump comes
from repair.

\begin{table}[t]
\centering
\small
\begin{tabular}{llccc}
\toprule
$K$ & Policy & Parse & Equiv. & Equiv./Parse \\
\midrule
1 & Greedy first & 0.924 & 0.368 & 0.399 \\
1 & Random parseable & 0.924 & 0.368 & 0.399 \\
1 & Verifier-ranked & 0.924 & 0.368 & 0.399 \\
1 & \vcsr repair-augmented & 0.924 & 0.368 & 0.399 \\
\midrule
4 & Greedy first & 0.924 & 0.368 & 0.399 \\
4 & Random parseable & 0.996 & 0.376 & 0.378 \\
4 & Verifier-ranked & 0.996 & 0.392 & 0.394 \\
4 & \vcsr repair-augmented & 0.996 & 0.392 & 0.394 \\
\midrule
8 & Greedy first & 0.924 & 0.368 & 0.399 \\
8 & Random parseable & 1.000 & 0.376 & 0.376 \\
8 & Verifier-ranked & 1.000 & 0.420 & 0.420 \\
8 & \vcsr repair-augmented & 1.000 & \textbf{0.772} & \textbf{0.772} \\
\bottomrule
\end{tabular}
\caption{Final fresh-seed aggregate metrics on seeds 51--55. Semantic
equivalence is the primary metric. Repair is applied only at the main $K=8$
operating point.}
\label{tab:main}
\end{table}

\subsection{Candidate-Pool Oracle}

Table~\ref{tab:oracle} reports candidate-pool diagnostics. At $K=8$, the
original generated pool contains at least one semantically equivalent candidate
for only 46.4\% of rows. This explains why verifier-only search has limited
headroom: even a perfect selector over the original pool could not solve rows
where no candidate is equivalent. Repair-augmented \vcsr reaches 77.2\% because
it edits the verifier-selected candidate and can create a new equivalent \pddl
problem. Repair exceeding oracle best-of-$K$ is therefore expected: repair adds
candidate quality rather than merely selecting better from the same pool.

\begin{table}[t]
\centering
\small
\begin{tabular}{cccc}
\toprule
$K$ & Avg. parseable candidates & Avg. equivalent candidates & Oracle best-of-$K$ \\
\midrule
1 & 0.924 & 0.368 & 0.368 \\
4 & 3.812 & 1.464 & 0.428 \\
8 & 7.564 & 2.888 & 0.464 \\
\bottomrule
\end{tabular}
\caption{Candidate-pool diagnostics on final seeds 51--55. Oracle best-of-$K$
uses gold equivalence labels and is only an offline upper bound for selecting
among the original generated candidates. Repair can exceed this diagnostic
because it creates a new candidate.}
\label{tab:oracle}
\end{table}

\subsection{Per-Seed Stability}

Table~\ref{tab:perseed} shows that repair improves every final seed. The
smallest gain is +0.28 and the largest gain is +0.42. Repair parse rate is high
across seeds, with a mean of 0.984.

\begin{table}[t]
\centering
\small
\begin{tabular}{ccccccc}
\toprule
Seed & Verifier $K=8$ & Repair $K=8$ & Delta & Helped & Hurt & Repair parse \\
\midrule
51 & 0.360 & 0.780 & +0.420 & 24 & 3 & 1.000 \\
52 & 0.420 & 0.780 & +0.360 & 21 & 3 & 0.980 \\
53 & 0.520 & 0.800 & +0.280 & 18 & 4 & 0.960 \\
54 & 0.340 & 0.760 & +0.420 & 22 & 1 & 0.980 \\
55 & 0.460 & 0.740 & +0.280 & 19 & 5 & 1.000 \\
\bottomrule
\end{tabular}
\caption{Per-seed final gate outcomes. Repair improves semantic equivalence on
all five untouched final seeds.}
\label{tab:perseed}
\end{table}

\subsection{Domain and Style Breakdown}

Table~\ref{tab:slices} highlights the main caveat. Repair is transformative on
\texttt{gripper}: verifier-ranked search solves only 2 of 124 rows, while repair
solves 98 of 124. In \texttt{blocksworld}, verifier-ranked search is already
stronger, and unconditional repair hurts some already-correct candidates:
blocksworld drops from 103/126 to 95/126. Overall, the gripper gains dominate,
but the domain asymmetry is important.

\begin{table}[t]
\centering
\small
\begin{tabular}{llccc}
\toprule
Slice & Policy & Equiv. count & Equiv. rate & Parse rate \\
\midrule
domain=blocksworld & Greedy first & 92/126 & 0.730 & 0.929 \\
domain=blocksworld & Verifier-ranked & 103/126 & 0.817 & 1.000 \\
domain=blocksworld & \vcsr repair-augmented & 95/126 & 0.754 & 1.000 \\
\midrule
domain=gripper & Greedy first & 0/124 & 0.000 & 0.919 \\
domain=gripper & Verifier-ranked & 2/124 & 0.016 & 1.000 \\
domain=gripper & \vcsr repair-augmented & 98/124 & 0.790 & 1.000 \\
\midrule
style=abstract/abstract & Greedy first & 28/126 & 0.222 & 0.849 \\
style=abstract/abstract & Verifier-ranked & 40/126 & 0.317 & 1.000 \\
style=abstract/abstract & \vcsr repair-augmented & 84/126 & 0.667 & 1.000 \\
\midrule
style=explicit/explicit & Greedy first & 64/124 & 0.516 & 1.000 \\
style=explicit/explicit & Verifier-ranked & 65/124 & 0.524 & 1.000 \\
style=explicit/explicit & \vcsr repair-augmented & 109/124 & 0.879 & 1.000 \\
\bottomrule
\end{tabular}
\caption{$K=8$ domain and style slices. Repair strongly improves gripper and
both style groups, but can hurt already-correct blocksworld selections.}
\label{tab:slices}
\end{table}

\subsection{Repair Outcomes}

Across the final gate, repair helped 104 rows, hurt 16 rows, and tied 130 rows.
The hurt rows are concentrated in blocksworld, where the verifier-selected
candidate is often already correct. Figure~\ref{fig:repairoutcomes} summarizes
the outcome counts.

\begin{figure}[t]
\centering
\fbox{
\begin{minipage}{0.65\linewidth}
\centering
\begin{tabular}{lr}
\toprule
Repair outcome & Count \\
\midrule
Helped & 104 \\
Hurt & 16 \\
Both success & 89 \\
Both fail & 41 \\
\bottomrule
\end{tabular}
\end{minipage}}
\caption{Final repair outcomes on seeds 51--55. Repair is a large net win but
not uniformly harmless.}
\label{fig:repairoutcomes}
\end{figure}

\section{Ablations and Analysis}

Figure~\ref{fig:funnel} summarizes the development evidence. The project began
with verifier training, but the final result is a system-level discovery:
verifier-only search helps, simple planner/search policies do not reliably
improve semantic selection, and repair provides the decisive jump.

\begin{figure}[t]
\centering
\fbox{
\begin{minipage}{0.86\linewidth}
\centering
Verifier training
$\rightarrow$ frozen round-4 verifier
$\rightarrow$ verifier-ranked $K=8$: 0.420
$\rightarrow$ solvability/search ablations: no reliable gain
$\rightarrow$ domain-aware repair
$\rightarrow$ repair-augmented $K=8$: 0.772.
\end{minipage}}
\caption{Evidence funnel from verifier-only search to repair-augmented \vcsr.}
\label{fig:funnel}
\end{figure}

\paragraph{Verifier-only search is useful but insufficient.}
Verifier-ranked selection improves over prompt-only generation at $K=8$,
raising semantic equivalence from 0.368 to 0.420 on the final gate. This
confirms that the verifier has useful ranking signal. However, the gain is
small compared with repair, so the verifier is best viewed as a scaffold for
identifying a candidate to repair rather than as the full solution.

\paragraph{Planner and solvability signals are weak semantic selectors.}
Planner and solvability policies did not pass the cached replay gate. This is
consistent with the motivation from Planetarium: solvability is not semantic
faithfulness. A solvable \pddl problem can still encode the wrong initial state
or goal.

\paragraph{Additional verifier training did not solve selection.}
We trained several verifier successors, including ranking-aligned and focused
pointwise variants. Some improved offline metrics or cached replay, but none
beat the frozen round-4 verifier sufficiently to become the paper-facing
selector. The remaining bottleneck was not solved by simply training the same
verifier harder.

\paragraph{Guarded repair did not fix blocksworld hurts.}
A simple verifier-score guard was tested on fresh seeds 67--71. It replicated
the large repair gain but did not reduce blocksworld hurts relative to
unconditional repair, so it is not promoted. This suggests that future repair
acceptance checks should use stronger structural signals rather than only the
same verifier score.

\section{Limitations}

The final claim is intentionally scoped. First, we evaluate only in-domain
Planetarium \texttt{blocksworld} and \texttt{gripper}; we make no final claim
for \texttt{floor-tile} or arbitrary new planning domains. Second, the repair
stage is domain-aware. This is appropriate for our in-domain system, but it may
limit immediate transfer to new domains. Third, repair is not uniformly
harmless: the final system is a large net win, but blocksworld shows that
repairing an already-correct candidate can introduce errors. Fourth, the
verifier remains imperfect and text-only; richer semantic validators,
graph-based verifiers, or structural repair acceptance checks may reduce the
remaining hurts. Finally, calibration and abstention were explored during
development but are not the final paper-facing mechanism.

\section{Conclusion}

\vcsr addresses a central failure mode in text-to-\pddl generation: generated
planning problems can be formally plausible yet semantically wrong. Our results
show that verifier-ranked best-of-$K$ search alone is helpful but insufficient.
The decisive improvement comes from adding one repair step to the
verifier-selected candidate. On untouched final seeds, repair-augmented \vcsr
raises semantic equivalence from 36.8\% for prompt-only generation to 77.2\%
at the main $K=8$ operating point, while maintaining high parseability. The
lesson is practical: if the goal is faithful planning formalization,
solvability and one-shot prompting are not enough; semantic search and targeted
repair are complementary ingredients. While our experiments are scoped to
Planetarium text-to-\pddl tasks, the underlying pattern is broader: systems that
translate natural language into structured executable artifacts need mechanisms
for checking and repairing meaning, not only syntax.

\appendix

\section{Claude-Family Robustness Benchmark}

After freezing the main result, we ran a small Claude-family robustness
benchmark on new seeds 72--74 with 10 rows per seed and model. This benchmark is
supporting evidence, not a replacement for the final seed 51--55 result.
Table~\ref{tab:modelbench} compares prompt-only generation against full \vcsr
for Claude Haiku 4.5, Claude Sonnet 4.5, and Claude Opus 4.6. All three improve
substantially under \vcsr. With prompt-only Opus, about 36.7\% of rows were
semantically correct; with full \vcsr around Opus, 90.0\% were semantically
correct, a +53.3 percentage-point gain on this small benchmark.

\begin{table}[t]
\centering
\small
\begin{tabular}{lccccc}
\toprule
Model & Prompt $K=1$ & \vcsr repair $K=8$ & Delta vs prompt & Random $K=8$ & Verifier $K=8$ \\
\midrule
Claude Haiku 4.5 & 0.400 & 0.900 & +0.500 & 0.300 & 0.467 \\
Claude Sonnet 4.5 & 0.500 & 0.933 & +0.433 & 0.467 & 0.533 \\
Claude Opus 4.6 & 0.367 & 0.900 & +0.533 & 0.433 & 0.433 \\
\bottomrule
\end{tabular}
\caption{Small post-freeze Claude-family robustness benchmark. These results
support the view that \vcsr is a useful wrapper, but the main paper evidence
remains the untouched final seed 51--55 gate.}
\label{tab:modelbench}
\end{table}

\section{Development Evidence Summary}

Round 4 remains the frozen verifier because it was the strongest checkpoint
under the replay and fresh held-out gates. Pairwise round 5, conservative
ranking round 6, and focused pointwise round 7 were useful diagnostics but were
not promoted. This history is important because it shows that the final gain is
not simply the result of training the verifier longer; the decisive improvement
comes from adding repair to the verifier-selected candidate.

\bibliographystyle{plainnat}
\bibliography{../references}

\end{document}
